%preamble: defining document parameters and packages

\documentclass[11pt]{report}         %defining documentclass
\usepackage[brazil] {babel}          %defining document language
\usepackage[utf8]{inputenc}          %defining to support portuguese a accentuation
\usepackage{authblk}                 %defining packages to better author format
\usepackage{indentfirst}             %defining packages to indention first paragraph's line
\usepackage[margin=2.5cm]{geometry}  %defining margin
\usepackage[table]{xcolor}           %defining table colors
\usepackage{pdfpages}                %enabling pdf attachment

%cover (title, author and date)

\title{\textbf{Preparatório para T2}}
\author{Arthur Siqueira Paz Teixeira \\ \& \\ Yure da Silva Lacerda de Barros} 
\affil{Disciplina: Circuitos Lógicos -- Professor: Fabio Marujo}
\date{\today}


\begin{document}                  %start document

\maketitle                        %generate the cover
\tableofcontents                  %generate the summary

\chapter {Introdução}             %First chapter

Como forma de otimizar os custos de produção, tempo de desenvolvimento e espaço ocupado pelo \textit{hardware} é de praxe que qualquer circuito
passe por um trabalho de simplificação usando técnicas tradicionais de simplificação. As principais fomas de simplificações incluem as tecnias da 
textbf{algebra booleana} e textbf{mapa de karnaugh}.

O presente documento referece as atividades desenvolvidas para a Aula Prática 2 (AP2) da disciplina de Circuitos Lógicos (EEL280). O objetivo 
central desta prática é o projeto, a simplificação e a implementação de circuitos combinacionais voltados para o acionamento de displays de 7 segmentos.

\chapter {Exercicios Propostos}

\includepdf[pages=-, pagecommand={\thispagestyle{plain}}]{/home/sparthur/git/logic-circuits-design/AP2/AP2-exercice-guide.pdf}

\chapter {Simplificação de Circuitos Simples}

\includepdf[pages=-, pagecommand={\thispagestyle{plain}}]{/home/sparthur/git/logic-circuits-design/AP2/Exercicio-1-resolvido.pdf}

\chapter {Display de 7 Segmentos}

\includepdf[pages=-, pagecommand={\thispagestyle{plain}}]{/home/sparthur/git/logic-circuits-design/AP2/display-de-7-segmentos.pdf}

\chapter {Montagem do Circuito}

Nesta ultima etapa é necessario o uso dos CI's 7404,7408 e 7432 para a lógica do circuito, de resistores para controle de tensão e proteção dos componentes, 
leds, para indicações visuais dos processos e estado do circuito, fios condutores para a corrente, o display de 7 segmentos (catodo comum) e um protoboard, 
onde o circuito vai ser montado e testado. Os componetes serão montados conforme demonstrado no capitulo anterior e posteriormente os testes serão conduzidos 
para a produção do relatório final onde maiores conclusões serão tiradas.

\end{document}
